The proposed algorithm alternates between estimating cluster memberships, updating GMM parameters, and refining the missing entries. It can be interpreted as an EM-like procedure nested inside a block-coordinate ascent routine.

\subsection{Responsibility Update}
Given the current completed matrix $\mat{X}$ and current parameters $\set*{\scal{\alpha}_k,\vect{\mu}_k,\mat{\Sigma}_k}_{k=1}^{\scal{K}}$, define the posterior responsibility of component $\scal{k}$ for sample $\vect{x}_j$ by
\[
    \scal{\gamma}_{jk}
    =
    p\paren*{\scal{z}_j = \scal{k} \given \vect{x}_j}
    =
    \frac{
        \scal{\alpha}_k \N\paren*{\vect{x}_j \given \vect{\mu}_k, \mat{\Sigma}_k}
    }{
        \sum_{\ell=1}^{\scal{K}}
        \scal{\alpha}_{\ell} \N\paren*{\vect{x}_j \given \vect{\mu}_{\ell}, \mat{\Sigma}_{\ell}}
    }.
\]
This is the standard E-step of model-based clustering~\cite{bouman1997cluster,melnykov2010finite}. The difference from the classical case is that $\vect{x}_j$ itself contains values generated by the previous completion step.

\subsection{Parameter Update with Fixed Completed Data}
Once $\scal{\gamma}_{jk}$ is available, the mixture coefficients, means, and covariance matrices are updated exactly as in a conventional GMM fitted on the current completed matrix:
\[
    \scal{N}_k = \sum_{j=1}^{\scal{n}} \scal{\gamma}_{jk},
    \qquad
    \scal{\alpha}_k = \frac{\scal{N}_k}{\scal{n}},
\]
\[
    \vect{\mu}_k
    =
    \frac{1}{\scal{N}_k}
    \sum_{j=1}^{\scal{n}}
    \scal{\gamma}_{jk}\vect{x}_j,
\]
\[
    \mat{\Sigma}_k
    =
    \frac{1}{\scal{N}_k}
    \sum_{j=1}^{\scal{n}}
    \scal{\gamma}_{jk}
    \paren*{\vect{x}_j - \vect{\mu}_k}
    \paren*{\vect{x}_j - \vect{\mu}_k}\trans.
\]

These equations reveal an important point often glossed over in short summaries: once the missing entries are temporarily filled, the model does not use a special incomplete-data likelihood at this stage. It uses the ordinary complete-data GMM update on the current surrogate matrix. The novelty lies in how that surrogate matrix is refreshed.

\subsection{Missing-Data Update with Fixed GMM Parameters}
The most distinctive step of the method is the update of $\vect{x}_{j,m}$. With all GMM parameters fixed, each sample solves
\[
    \argmax_{\vect{x}_{j,m}}
    \sum_{k=1}^{\scal{K}}
    \scal{\alpha}_k
    \N\paren*{\vect{x}_j \given \vect{\mu}_k, \mat{\Sigma}_k}
    \quad
    \text{s.t. }
    \vect{x}_{j,o} = \vect{x}^{\,\text{obs}}_{j,o}.
\]

To make the derivation explicit, partition each mean vector and precision matrix according to observed and missing coordinates:
\[
    \vect{\mu}_k =
    \begin{bmatrix}
        \vect{\mu}_{k,o}\\
        \vect{\mu}_{k,m}
    \end{bmatrix},
    \qquad
    \mat{\Sigma}_k\inv =
    \begin{bmatrix}
        \mat{\Lambda}_{k,oo} & \mat{\Lambda}_{k,om}\\
        \mat{\Lambda}_{k,mo} & \mat{\Lambda}_{k,mm}
    \end{bmatrix}.
\]
The update derived in the paper is
\begin{equation}
\begin{aligned}
    \vect{x}_{j,m}
    &=
    \paren*{
        \sum_{k=1}^{\scal{K}}
        \scal{P}_{jk}\mat{\Lambda}_{k,mm}
    }\inv \\
    &\quad \times
    \sum_{k=1}^{\scal{K}}
    \scal{P}_{jk}
    \paren*{
        \mat{\Lambda}_{k,mo}\vect{\mu}_{k,o}
        +
        \mat{\Lambda}_{k,mm}\vect{\mu}_{k,m}
        -
        \mat{\Lambda}_{k,mo}\vect{x}_{j,o}
    },
\end{aligned}
\end{equation}
where
\[
    \scal{P}_{jk} = \scal{\alpha}_k \N\paren*{\vect{x}_j \given \vect{\mu}_k, \mat{\Sigma}_k}.
\]

This equation is the mathematical core of the method. It says that the missing block is reconstructed from a weighted combination of component-wise conditional geometry. The weights depend on current mixture plausibility, and the linear system is expressed in the \highlight{precision domain} rather than directly in covariance form. That is why the update resembles a compromise across local quadratic energies induced by each Gaussian component.

\subsection{Interpretation as Conditional Gaussian Reconstruction}
For a single Gaussian component, the optimal missing-value estimate under a quadratic log-density would reduce to a conditional mean expression driven by the relation between observed and missing coordinates. The proposed method generalizes that logic to a mixture model. Instead of selecting one component, it interpolates across components through the weights $\scal{P}_{jk}$. This yields a softer and more adaptive estimator than hard assignment to a single cluster center.

\begin{table}[t]
    \caption{Two-stage imputation versus the one-stage GMM framework}
    \label{tab:one-stage-vs-two-stage}
    \centering
    \begin{tabularx}{\linewidth}{@{}lX X@{}}
        \toprule
        Aspect & Two-stage pipeline & Proposed one-stage method \\
        \midrule
        Imputation target & Generic feature reconstruction & Cluster-aware likelihood improvement \\
        Feedback from clustering & None & Directly affects missing-value updates \\
        Use of covariance geometry & Often weak or absent & Explicit through $\mat{\Sigma}_k$ and $\mat{\Sigma}_k\inv$ \\
        Robustness at high missingness & Tends to deteriorate quickly & Usually more stable if initialization is adequate \\
        Optimization view & Sequential heuristics & Alternating ascent on a joint objective \\
        \bottomrule
    \end{tabularx}
\end{table}

\subsection{Algorithmic Interpretation}
The whole procedure can be summarized as follows:
\begin{enumerate}
    \item initialize missing entries with a simple imputer such as column means;
    \item initialize GMM parameters on the completed matrix;
    \item compute posterior responsibilities $\scal{\gamma}_{jk}$;
    \item update $\scal{\alpha}_k$, $\vect{\mu}_k$, and $\mat{\Sigma}_k$;
    \item update each $\vect{x}_{j,m}$ while keeping $\vect{x}_{j,o}$ unchanged;
    \item repeat until the log-likelihood improvement falls below a tolerance.
\end{enumerate}

This is best viewed as block-coordinate ascent on a non-convex objective. Each block update is designed to avoid decreasing the current objective, so the procedure inherits a monotonic improvement property. However, because the optimization is non-convex, the destination is a local optimum rather than a globally optimal solution.

\subsection{Complexity and Numerical Stability}
The paper reports a complexity of $\mathcal{O}\paren*{\scal{t}\scal{K}\scal{n}\scal{d}^3}$, with $\scal{t}$ iterations, $\scal{K}$ clusters, $\scal{n}$ samples, and $\scal{d}$ dimensions. The cubic term comes from matrix inversion in the data-completion step. In practice, this immediately implies three engineering concerns:
\begin{itemize}
    \item covariance regularization is mandatory when clusters are small or features are highly correlated;
    \item poor initialization can trap the method in unstable local optima;
    \item the method becomes expensive in high dimensions unless one constrains $\mat{\Sigma}_k$, for example to diagonal, tied, or low-rank structures.
\end{itemize}

These issues are not implementation details; they shape the actual usability of the method.

\subsection{Theoretical Boundary of the Current Formulation}
The method remains elegant, but it is not equivalent to exact maximum likelihood under incomplete observations in the classical missing-data EM sense. A textbook incomplete-data EM would integrate over missing entries in the E-step, whereas this method updates them by deterministic optimization. That difference is not a flaw by itself, but it must be stated clearly because it affects how we interpret the algorithm's statistical guarantees.
